\begin{figure}
\begin{center}

\setlength{\unitlength}{0.001in}

\begin{picture}(3000,1900)(0,-1600)

\multiput(0300,-1400)(0,300){6}{\circle {50}}
\multiput(0600,-1400)(0,300){6}{\circle {50}}
\multiput(0900,-1400)(0,300){6}{\circle {50}}
\multiput(1200,-1400)(0,300){6}{\circle {50}}
\multiput(1500,-1400)(0,300){6}{\circle {50}}
\multiput(1800,-1400)(0,300){6}{\circle {50}}
\multiput(2100,-1400)(0,300){6}{\circle {50}}
\multiput(2400,-1400)(0,300){6}{\circle {50}}
\multiput(2700,-1400)(0,300){6}{\circle {50}}

\multiput(0450,-1550)(300,0){8}{\dashbox{50}(0000,1800){}}
\multiput(0150,-1250)(0,300){5}{\dashbox{50}(2700,0000){}}

\multiput(0400,-1450)(0,800){3}{\circle*{100}}
\multiput(1400,-1450)(0,800){3}{\circle*{100}}
\multiput(2400,-1450)(0,800){3}{\circle*{100}}

\put(895,-1055){\framebox(1010,810){}}
\put(900,-1050){\framebox(1000,800){}}
\put(905,-1045){\framebox(990,790){}}

\end{picture}
\end{center}

\caption{双向嵌套方法中协调低等级网格与高等级网格的概念示意。$\circ$ 和虚线分别表示高等级网格点和网格单元，$\bullet$ 和实线分别表示低等级网格及其中心网格单元。}
\label{fig:nest2}
\botline
\end{figure}
